\hypertarget{appendix-ab31ef447dbf}{}\label{referencestart:appendix-ab31ef447dbf}
\section{Matrix multiplication and composition order}\label{appendix-ab31ef447dbf-matrix-multiplication-and-composition-order}

\textbf{Curriculum depth:} Introduced · \textbf{First seen:} Phase7 · \textbf{Latest development:} Phase 9

\subsection{Definition and mechanism}\label{appendix-ab31ef447dbf-definition-and-mechanism}

The product AB applies B first for column-vector actions. Its entries sum A(i,k)B(k,j) over the shared dimension. Composition is associative over exact scalars and usually noncommutative; floating-point evaluation orders need not agree bitwise.

\textbf{Reconstruction cue:} Rotation followed by translation differs from translation followed by rotation.

\subsection{Prerequisites}\label{appendix-ab31ef447dbf-prerequisites}

\begin{itemize}
\tightlist
\item
  \hyperref[appendix-bd39c3bf5d64]{Matrices as linear maps and basis images}
\end{itemize}

\subsection{Used by}\label{appendix-ab31ef447dbf-used-by}

\begin{itemize}
\item
  \hyperref[appendix-175ca5ff874d]{Matrix transpose symmetry and trace}
\item
  \hyperref[appendix-5f723fbd0593]{Euler angles and gimbal singularity}
\item
  \hyperref[appendix-0b3be18411f9]{Eigenvalues and eigenspaces} --- Extends this foundation into linear-systems and decomposition analysis.
\item
  \hyperref[appendix-1fad5e7f36fa]{Singular value decomposition} --- Extends this foundation into linear-systems and decomposition analysis.
\item
  \hyperref[appendix-0f4bc2bf2b1f]{Matrix functions} --- Extends this foundation into linear-systems and decomposition analysis.
\item
  \hyperref[appendix-f6e77f4a200f]{Gaussian elimination and pivoting} --- Extends this foundation into the Phase 8 dependency spine.
\item
  \hyperref[appendix-40df0855b4a7]{Automatic differentiation and backpropagation}
\item
  \hyperref[appendix-d8a7ddb675eb]{Jacobians, pushforwards, and pullbacks}
\end{itemize}

\subsection{Recall and next depth}\label{appendix-ab31ef447dbf-recall-and-next-depth}

Reconstruct the example, state its conventions and find an input that violates its preconditions. Deeper algorithms and measured production performance remain separate from the present model.

\subsection{Source and navigation}\label{appendix-ab31ef447dbf-source-and-navigation}

\href{../02-Mathematics-for-Visual-Computing/main.pdf\#nameddest=01-foundations-i-ch018-keep-the-promises-visible-while-moving-between-representations}{Continuous explanation (Mathematics for Visual Computing, Mathematics for Visual Computing, Chapter 2)} · \hyperref[appendix-fd090d982a15-phase-7---parts-601-700]{Coverage and source gap} · \hyperref[appendix-a103e7846d21-phase-7---source-analysis-and-verification]{Sources and verification}

\subsection{Full teaching explanation}\label{appendix-ab31ef447dbf-full-teaching-explanation}

\begin{itemize}
\tightlist
\item
  \href{../02-Mathematics-for-Visual-Computing/main.pdf\#nameddest=02-foundations-ii-ch021-jacobians-curvature-and-local-approximation}{Phase 9 teaching section (Mathematics for Visual Computing, Mathematics for Visual Computing, Chapter 5)}
\item
  \hyperref[appendix-adf48f5dc5c1]{Visual Recall/09 - Calculus and Optimization - Visual Recall.md} · Phase 9 visual recall
\end{itemize}
